%% main.tex, Paper 2: Spectral Surrogates for Rough Stochastic Volatility
%%
%% Working title: Distributional Learning of Rough Volatility via Spectral
%%        Surrogates: a Cramér–Bochner PINN for fBm-Driven SDEs
%%
%% Target: SIAM Journal on Financial Mathematics / Quantitative Finance
%%
%% Build: pdflatex main → bibtex main → pdflatex main ×2
%% ─────────────────────────────────────────────────────────────────────────────

\documentclass[11pt,a4paper]{article}

\usepackage{amsmath, amssymb, amsthm}
\usepackage{mathtools}
\usepackage[margin=2.5cm]{geometry}
\usepackage{hyperref}
\usepackage{cleveref}
\usepackage{booktabs}
\usepackage{enumitem}
\usepackage{graphicx}
\graphicspath{{../02_rough_vol/figures/}}
\usepackage{xcolor}
\usepackage{microtype}
\usepackage{natbib}

% ── Theorem environments ──────────────────────────────────────────────────────
\newtheorem{theorem}{Theorem}[section]
\newtheorem{lemma}[theorem]{Lemma}
\newtheorem{corollary}[theorem]{Corollary}
\newtheorem{proposition}[theorem]{Proposition}
\newtheorem{remark}[theorem]{Remark}
\newtheorem{assumption}{Assumption}
\renewcommand{\theassumption}{A\arabic{assumption}}

% ── Notation ─────────────────────────────────────────────────────────────────
\newcommand{\R}{\mathbb{R}}
\newcommand{\E}{\mathbb{E}}
\newcommand{\Prob}{\mathbb{P}}
\newcommand{\Law}{\mathcal{L}}
\newcommand{\W}{W_2}        % Wasserstein-2
\newcommand{\Hcal}{\mathcal{H}}
\newcommand{\Lcal}{\mathcal{L}}
\newcommand{\norm}[1]{\|#1\|}
\newcommand{\HS}{\mathrm{HS}}
\newcommand{\dt}{\,\mathrm{d}t}
\newcommand{\dw}{\,\mathrm{d}W}
\newcommand{\TODO}[1]{\textcolor{red}{[\textbf{TODO:} #1]}}

% ─────────────────────────────────────────────────────────────────────────────
\begin{document}

\title{Distributional Learning of Rough Volatility via Spectral Surrogates:\\
    a Cramér PINN for Fractional Brownian Motion--Driven SDEs}

\author{Hubert Lipski}
\date{\today}
\maketitle

% ─────────────────────────────────────────────────────────────────────────────
\begin{abstract}
Rough stochastic volatility models, in which the log-variance is driven by a
fractional Brownian motion with Hurst exponent $H \approx 0.1$, reproduce
key empirical features of equity implied volatility but require expensive
quasi-Monte Carlo simulation for option pricing. We introduce the
\emph{Cramér PINN}, a neural-network surrogate trained to match the
power-spectral density (PSD) of fractional Gaussian noise (fGn), the
stationary increment process of fBm. The surrogate requires no
numerical solver at inference time and generates variance paths by a
single forward pass. We prove a $W_2$ convergence theorem showing that
matching the exact fGn PSD forces the surrogate's law to converge to the
reference process, with an explicit Hurst-exponent consistency corollary:
if the training loss is $\varepsilon$, then $|\hat{H} - H| \leq C\sqrt{\varepsilon}$.
On the rough Bergomi model ($H=0.1$, $\eta=1.9$, $\xi_0=0.04$), a linear
spectral filter surrogate trained on the exact fGn PSD achieves an
implied-volatility surface RMSE of \textbf{1.33~vol~points} vs.\ the MC
reference, compared to \textbf{2.27~vp} for the asymptotic
$C|f|^{1-2H}$ approximation. A nonlinear MLP surrogate trained with
the same spectral loss reaches \textbf{1.21~vp}, confirming that the loss
is architecture-agnostic. Under rough Heston with $\rho=-0.9$, we show that any surrogate matching
only the marginal distribution of $V$ carries an \textbf{irreducible ATM
delta bias of $-0.27$} ($35\%$ relative). We further show that circular
FFT surrogates fail catastrophically at $\rho\neq 0$ (204~vp, delta $>1$)
due to non-causal wrap-around, and that a \textbf{CausalSpectralFilter}
(lower-triangular Toeplitz convolution, same PSD loss) resolves this,
achieving \textbf{3.26~vp} vs.\ the BLP reference and reducing ATM delta
bias to $-0.07$. Training the causal filter on the reference's own
RL-fBm PSD eliminates the residual entirely: \textbf{2.57~vp} and ATM
delta bias $\lesssim 0.01$ (Monte-Carlo-noise level). On rough Heston,
the same spectral causal architecture cuts the cross-model bias from
$+0.49$ to $-0.15$; a Volterra-aware autoregressive surrogate with
$\sqrt{V_j}$ diffusion closes the residual to Monte-Carlo noise
(\textbf{0.80~vp}, bias $\leq 0.005$). The same architecture, trained
end-to-end from a misspecified init, calibrates to a live SPY option
chain across three expirations ($32$, $88$, $179$ days; $26$ strikes
total) at \textbf{$3.97$~vp overall IV RMSE} in $3$~minutes on a
single GPU, with shared parameters across maturities.
\end{abstract}

\tableofcontents
\newpage

% ═════════════════════════════════════════════════════════════════════════════
\section{Introduction}
\label{sec:intro}
% ═════════════════════════════════════════════════════════════════════════════

The empirical observation that log-volatility behaves like a fractional
Brownian motion with Hurst exponent $H \approx 0.1$
\citep{gatheral2018volatility} has led to a family of \emph{rough
volatility} models that fit the implied volatility surface and the term
structure of variance swaps more accurately than classical stochastic
volatility models \citep{bayer2016pricing,ElEuch2019}. The prototypical
example is the rough Bergomi (rBergomi) model \citep{bayer2016pricing}:
\begin{equation}
 V_t \;=\; \xi_0 \exp\!\Bigl(\eta W^H_t - \tfrac{1}{2}\eta^2 t^{2H}\Bigr),
 \qquad
 \frac{dS_t}{S_t} \;=\; \sqrt{V_t}\,dZ_t,
 \label{eq:rbergomi}
\end{equation}
where $W^H$ is a fractional Brownian motion (fBm) with Hurst exponent
$H$, and $(Z_t, W^H_t)$ are correlated.

\paragraph{The computational bottleneck.}
Pricing options under \eqref{eq:rbergomi} requires Monte Carlo simulation
of fBm paths, which is inherently sequential and cannot exploit the
structure of the problem. Quasi-Monte Carlo methods improve convergence
rates \citep{bennedsen2017hybrid} but do not eliminate the cost. For
applications such as deep hedging \citep{buhler2019deep}, which require
millions of simulated paths, the cost is prohibitive.

\paragraph{This paper.}
We propose to replace the Monte Carlo simulator with a \emph{spectral
surrogate}: a neural network trained to generate paths whose
power-spectral density (PSD) matches the exact PSD of fractional Gaussian
noise (fGn), the stationary increment process of fBm. This is a direct
extension of the Bochner PINN framework of \citet{TODO-paper1}
from stationary processes (where Bochner's theorem applies) to processes
with stationary increments (where the Cramér spectral representation
applies).

\paragraph{Contributions.}
\begin{enumerate}
 \item \textbf{Cramér PINN architecture.} A linear spectral filter
  (SpectralFilter, 126 parameters) and a nonlinear MLP (BochnerMLP,
  41k parameters) both trained to minimise the log-PSD MSE against the
  exact fGn PSD. Smart initialisation from the analytical optimum
  reduces training to fine-tuning.

 \item \textbf{Convergence theorem.} Under mild regularity assumptions,
  matching the fGn PSD in $L^2$ forces $W_2$-convergence of the
  surrogate to the reference process. Corollary: $|\hat{H} - H| \leq
  C\sqrt{\varepsilon}$.

 \item \textbf{Financial validation.} On the rBergomi model, the
  spectral surrogate achieves 1.21--1.33~vol-point IV surface RMSE
  vs.\ MC truth, near the Monte Carlo noise floor with 10k paths.

 \item \textbf{Exact vs.\ asymptotic PSD.} The asymptotic power-law
  $C|f|^{1-2H}$ is a 1.7$\times$ worse spectral target than the
  exact Monte Carlo fGn PSD across all maturities, most severely at
  long maturities ($T=1$: 3.20 vs.\ 1.89~vp).

 \item \textbf{Leverage-effect, Greek bias, and causal architecture.}
  Under rough Heston ($\rho=-0.9$), any marginal-$V$ surrogate carries
  an irreducible ATM delta bias of $-0.27$ ($35\%$). A circular
  $z$-Cholesky attempt fails (204~vp, delta $>1$), with non-causal FFT
  wrap-around as the precise obstruction. A \textbf{CausalSpectralFilter}
  (lower-triangular Toeplitz, same PSD loss) resolves this, achieving
  \textbf{3.26~vp} vs.\ BLP rBergomi reference (63$\times$ improvement)
  and reducing ATM delta bias to $-0.07$; matching the training PSD to
  the reference's RL-fBm distribution drops both to \textbf{2.57~vp}
  and $\lesssim 0.01$, Monte-Carlo-noise level. On rough Heston the
  same spectral architecture cuts the cross-model bias from $+0.49$
  to $-0.15$ (41$\times$ IV RMSE improvement, recovering $\sim\!42\%$
  of the leverage effect), and the remaining gap closes to MC noise
  (\textbf{0.80~vp}, bias $\lesssim 0.005$) with a Volterra-aware
  autoregressive surrogate, pinpointing the V-conditional diffusion
  ($\sqrt{V_j}$ in the noise term) as the precise architectural
  requirement.

 \item \textbf{Trainability and live-market validation.}
  The full parameter set $(\theta, c_{\rm drift}, c_{\rm stoch}, \beta,
  \{w_i\})$ is trainable end-to-end from a misspecified init via
  backprop through the autoregressive Volterra update; joint training
  on synthetic rough Heston with a price + V-moment loss reaches
  \textbf{1.85~vp} (close to the architectural $0.80$~vp ceiling).
  On a live SPY option chain (Yahoo Finance, three expirations spanning
  $32$ to $179$ days, $26$ strikes total across $K/S_0 \in [0.85, 1.15]$),
  the same surrogate calibrates simultaneously to all maturities at
  \textbf{$3.97$~vp overall IV RMSE} (short-end $2.27$~vp) in
  $3$~minutes on a GTX~950, within $2\times$ of the bid/ask-implied
  IV uncertainty at the short end. All parameters are \emph{shared}
  across maturities; the $\sim\!5$~vp long-end residual reflects the
  known forward-variance term-structure limitation of stationary
  rough-vol calibrations. A Markovian variant
  of the kernel ($M = 5$ exponentials) gives a $7\times$ inference
  speedup with no loss in IV-RMSE accuracy.
\end{enumerate}

\paragraph{Relation to Paper 1.} \citet{TODO-paper1} prove a $W_2$
convergence theorem for Bochner PINNs applied to Lévy-driven SPDEs whose
solutions are stationary. fBm is not stationary but has stationary
increments; the Cramér representation replaces Bochner's theorem and all
other proof steps carry over unchanged. \Cref{sec:theory} makes this
precise.

% ═════════════════════════════════════════════════════════════════════════════
\section{Background}
\label{sec:background}
% ═════════════════════════════════════════════════════════════════════════════

\subsection{The rough Bergomi model}

Let $(\Omega, \mathcal{F}, \mathbb{P})$ be a probability space.
The \emph{rough Bergomi} model \citep{bayer2016pricing} specifies the
stock price $S_t$ and instantaneous variance $V_t$ by
\eqref{eq:rbergomi}, where:
\begin{itemize}
 \item $W^H_t$ is a fractional Brownian motion with Hurst exponent
  $H \in (0, \tfrac{1}{2})$, defined via the Riemann--Liouville
  representation
  $W^H_t = \int_0^t (t-s)^{H-1/2} \dw_s / \Gamma(H + \tfrac{1}{2})$.
 \item $\xi_0 > 0$ is the initial forward variance (spot vol
  $\sqrt{\xi_0}$).
 \item $\eta > 0$ is the vol-of-vol parameter.
 \item The Itô correction $-\frac{1}{2}\eta^2 t^{2H}$ ensures
  $\E[V_t] = \xi_0$ for all $t$.
\end{itemize}

\paragraph{Why fBm?} The process $\log V_t = \log \xi_0 + \eta W^H_t -
\frac{1}{2}\eta^2 t^{2H}$ has increments driven by fBm. For $H < 1/2$,
$W^H$ has negatively correlated increments and paths that are
H\"older-continuous of order $H - \varepsilon$ but not differentiable.
This \emph{roughness} is precisely what the empirical evidence
demands \citep{gatheral2018volatility}.

\begin{figure}[t]
 \centering
 \includegraphics[width=\textwidth]{roughness_gallery}
 \caption{%
  \textbf{Roughness gallery.}
  Simulated fBm paths $W^H_t$ (top row) and instantaneous variance
  $V_t = \xi_0\exp(\eta W^H_t - \tfrac{1}{2}\eta^2 t^{2H})$ (bottom row)
  for $H \in \{0.50, 0.30, 0.10, 0.05\}$, with $\eta=1.9$, $\xi_0=0.04$.
  Decreasing $H$ produces increasingly rough, rapidly mean-reverting
  paths. The empirically calibrated value $H \approx 0.1$ (third
  column) generates the jagged, persistent variance clustering seen
  in equity markets. Standard Brownian motion ($H=0.5$, left) is
  manifestly too smooth.
 }
 \label{fig:roughness-gallery}
\end{figure}

\subsection{Fractional Gaussian noise and its PSD}
\label{sec:fgn-psd}

The \emph{fractional Gaussian noise} (fGn) is the increment process
$\Delta W^H_k = W^H_{k/N} - W^H_{(k-1)/N}$, $k = 1, \ldots, N$.
Unlike fBm, fGn is stationary. Its power spectral density is
\begin{equation}
 S_H(f) \;=\; \sigma^2 \int_{-\infty}^{\infty}
  e^{-2\pi i f x}\, r(x) \,\mathrm{d}x,
 \qquad
 r(n) = \tfrac{1}{2}(|n+1|^{2H} + |n-1|^{2H} - 2|n|^{2H}),
 \label{eq:fgn-psd}
\end{equation}
where $r(n)$ is the autocovariance at lag $n$. The \emph{asymptotic}
behaviour for $f \to 0$ is the power law
\begin{equation}
 S_H(f) \;\sim\; C_H |f|^{1-2H}, \qquad
 C_H = \Gamma(2H+1)\sin(\pi H)/(2\pi).
 \label{eq:power-law}
\end{equation}
For $H = 0.1$: $1 - 2H = 0.8$ (the PSD \emph{increases} with frequency,
reflecting roughness). For $H = 0.5$ (standard BM): $S_{1/2}(f) =
\text{const}$ (white noise).

\paragraph{Key point.} The asymptotic law \eqref{eq:power-law} holds
only as $f \to 0$. At mid-frequencies the exact PSD \eqref{eq:fgn-psd}
deviates from the power law, and this deviation drives the difference
between our two training targets (see \cref{sec:experiments}).

\subsection{Cramér spectral representation}

For a centred Gaussian process $u$ with stationary increments, the
\emph{Cramér representation} \citep{cramer1942harmonic} is
\begin{equation}
 u(t) - u(0) \;=\; \int_{-\infty}^{\infty}
  \bigl(e^{i\omega t} - 1\bigr) Z(\mathrm{d}\omega),
 \label{eq:cramer}
\end{equation}
where $Z$ is an orthogonal random measure with
$\E|Z(\mathrm{d}\omega)|^2 = S(\omega)\,\mathrm{d}\omega$.
This is the analogue of Bochner's theorem for processes with stationary
levels; it plays the same role in the Paper 2 proof that Bochner plays
in Paper 1.

% ═════════════════════════════════════════════════════════════════════════════
\section{Spectral Surrogates for fBm-Driven SDEs}
\label{sec:method}
% ═════════════════════════════════════════════════════════════════════════════

\subsection{Architecture}

We train a neural network $f_\theta : \R^N \to \R^N$ to generate fGn
increment sequences $\Delta W \approx f_\theta(z)$ from i.i.d.\ Gaussian
input $z \sim \mathcal{N}(0, I_N)$. Two architectures:

\paragraph{SpectralFilter.} A learnable per-frequency log-amplitude filter:
\begin{equation}
 f_\theta(z) \;=\; \mathrm{IRFFT}\!\bigl(A \cdot \mathrm{RFFT}(z)\bigr),
 \qquad
 A_k = \begin{cases} 0 & k = 0 \\ \exp(\alpha_k) & k = 1,\ldots, N/2 \end{cases}
 \label{eq:spectral-filter}
\end{equation}
where $\boldsymbol{\alpha} \in \R^{N/2}$ are the learnable parameters
(126 for $N=252$). The DC component is fixed to zero to enforce
zero-mean output and prevent unpenalised drift accumulation in cumulative
sums. Setting $A_k = 0$ at DC ensures $\E[\Delta W_k] = 0$ and
$\mathrm{Var}[\sum_k \Delta W_k]$ is determined entirely by the
spectral shape.

\textbf{Optimal solution.} The output $f_\theta(z)$ is Gaussian by
linearity. At the training optimum, $\E[\mathrm{psd}(f_\theta(z), f_k)]
= S_H(f_k)$ for each frequency $f_k$, which gives the analytic solution
\begin{equation}
 \alpha_k^* = \tfrac{1}{2}\log\!\bigl(S_H(f_k) \cdot N / (2\,\mathrm{d}t)\bigr).
 \label{eq:smart-init}
\end{equation}
We use \eqref{eq:smart-init} as the initialisation (\emph{smart init}),
reducing training to fine-tuning (initial loss $\approx 0.018$ vs.\
$\approx 3.1$ from zero-initialisation).

\paragraph{BochnerMLP.} A small fully-connected network with Tanh
activations:
\begin{equation}
 f_\theta(z) \;=\; \mathrm{MLP}(z) - \overline{\mathrm{MLP}(z)},
 \label{eq:mlp}
\end{equation}
where the mean subtraction enforces zero-mean output. Architecture:
$N \to 64 \to 64 \to 64 \to N$, 40\,892 parameters for $N=252$.
The MLP generates non-Gaussian paths; the spectral loss is
architecture-agnostic.

\subsection{Spectral loss}

Let $\hat{S}(\Delta W, f_k) = (2\,\mathrm{d}t / N)|[\mathrm{RFFT}(\Delta W
- \overline{\Delta W})]_k|^2$ be the one-sided periodogram. The
\emph{Bochner spectral loss} is the log-MSE:
\begin{equation}
 \mathcal{L}(\theta) \;=\;
 \frac{1}{N/2} \sum_{k=1}^{N/2}
 \Bigl(\log \hat{S}(f_\theta(z), f_k) - \log S_H(f_k)\Bigr)^2,
 \label{eq:loss}
\end{equation}
averaged over a mini-batch. The log transformation gives equal
relative weight to each frequency, preventing the high-power
low-frequency bins from dominating.

Two training targets:
\begin{enumerate}
 \item \textbf{Exact}: $S_H(f_k)$ estimated by averaging 100\,000
  Monte Carlo fGn periodograms (Wood--Chan simulation, cached to disk).
 \item \textbf{Asymptotic}: $S_H^{\mathrm{asymp}}(f_k) = C|f_k|^{1-2H}$
  normalised to the same total power as the exact target.
\end{enumerate}

\begin{figure}[t]
 \centering
 \includegraphics[width=\textwidth]{psd_comparison}
 \caption{%
  \textbf{Exact vs.\ asymptotic fGn spectral density, $H=0.1$.}
  \emph{Main panel:} log--log plot of the exact Monte Carlo PSD
  (blue, averaged over $10^5$ Wood--Chan paths) and the asymptotic
  power-law $C|f|^{1-2H}$ (orange, normalised to equal total power).
  The two targets agree in slope but diverge in amplitude across the
  full frequency range.
  \emph{Inset:} relative error $(S_{\mathrm{exact}}-S_{\mathrm{asymp}})/S_{\mathrm{exact}}$;
  the asymptotic target over-estimates power at mid-frequencies by
  up to 15\%, which biases long-maturity option prices
  (see Table~\ref{tab:iv-rmse} and \cref{fig:iv-surrogate}).
 }
 \label{fig:psd-comparison}
\end{figure}

\subsection{Variance path generation}

Given trained fGn increments $\Delta W = f_\theta(z)$, construct:
\begin{align}
 \hat{W}_t &= \sum_{k : t_k \leq t} \Delta W_k \quad \text{(fBm approximation)}, \\
 \hat{V}_t &= \xi_0 \exp\!\Bigl(\eta \hat{W}_t
  - \tfrac{1}{2}\eta^2 \widehat{\mathrm{Var}}[\hat{W}_t]\Bigr),
 \label{eq:vhat}
\end{align}
where $\widehat{\mathrm{Var}}[\hat{W}_t]$ is the cross-sectional
empirical variance of $\hat{W}_t$ across the batch. This ensures
$\E[\hat{V}_t] = \xi_0$ for any surrogate whose paths may not have
exactly the right total variance.

% ═════════════════════════════════════════════════════════════════════════════
\section{Convergence Theory}
\label{sec:theory}
% ═════════════════════════════════════════════════════════════════════════════

\subsection{Assumptions}

\begin{assumption}[Stationary increments]\label{ass:si}
 Both $u_\theta$ and $u_{\mathrm{ref}}$ are centred Gaussian processes
 on $[0,T]$ with stationary increments; they admit the Cramér
 representation \eqref{eq:cramer} with spectral densities
 $S_\theta$ and $S_{\mathrm{ref}}$, respectively.
\end{assumption}

\begin{assumption}[Spectral non-degeneracy]\label{ass:nd}
 There exist $0 < \omega_0 < \Omega < \infty$ and $\mu_{\min} > 0$
 such that $S_i(\omega) \geq \mu_{\min}$ for $\omega \in
 [\omega_0,\Omega]$, $i \in \{\theta, \mathrm{ref}\}$.
\end{assumption}

\begin{assumption}[Training tolerance]\label{ass:tol}
 The spectral loss satisfies $\mathcal{L}(\theta) \leq \varepsilon$.
\end{assumption}

\begin{assumption}[Gaussianity]\label{ass:gauss}
 $u_{\mathrm{ref}}$ is a Gaussian process (satisfied by fBm).
\end{assumption}

\subsection{Main theorem}

\begin{theorem}[Cramér PINN convergence]\label{thm:main}
 Under Assumptions~\ref{ass:si}--\ref{ass:gauss},
 \begin{equation}
  W_2\!\bigl(\mathcal{L}(u_\theta),\,\mathcal{L}(u_{\mathrm{ref}})\bigr)
  \;\leq\;
  \underbrace{\frac{C_T}{\mu_{\min}^{1/2}}
  \bigl\|S_\theta - S_{\mathrm{ref}}\bigr\|_{L^2([\omega_0,\Omega])}}_{\text{spectral term}}
  +
  \underbrace{C_2\,N_s^{-\beta_H}}_{\text{discretisation}}
  +
  \underbrace{C_3\,\omega_0^{-(1-2H)/2}}_{\text{cutoff}},
  \label{eq:main-bound}
 \end{equation}
 where $C_T = T/(2\sqrt{2\pi})$, $\beta_H = \frac{1}{2}\min(H - 1/2 + d/2, 1)$,
 and $N_s$ is the number of spectral collocation points.
\end{theorem}

\begin{proof}[Proof sketch]
 Step 1 (Gelbrich): Gaussian measures on $L^2([0,T])$ satisfy
 $W_2^2 = \|C_\theta^{1/2} - C_{\mathrm{ref}}^{1/2}\|_\HS^2$
 (Gelbrich formula \citep{gelbrich1990formula}).
 Step 2 (Cramér-to-Gelbrich, see Lemma~\ref{lem:cramer-gelbrich}):
 $W_2^2 \leq \mu_{\min}^{-1/2}\|C_\theta - C_{\mathrm{ref}}\|_\HS$.
 Step 3 (Parseval): $\|C_\theta - C_{\mathrm{ref}}\|_\HS^2 \leq
 T^2 \|S_\theta - S_{\mathrm{ref}}\|_{L^2}^2$.
 The discretisation and cutoff terms follow by standard quadrature
 bounds (cf.\ \citet{TODO-paper1}, Appendix~A).
\end{proof}

\subsection{Cramér-to-Gelbrich lemma}

\begin{lemma}[Cramér-to-Gelbrich]\label{lem:cramer-gelbrich}
 Under Assumptions~\ref{ass:si}--\ref{ass:nd},
 \[
  W_2\!\bigl(\mathcal{L}(u_1), \mathcal{L}(u_2)\bigr)
  \;\leq\;
  \frac{T}{2\sqrt{2\pi}\,\mu_{\min}^{1/2}}
  \bigl\|S_1 - S_2\bigr\|_{L^2([\omega_0,\Omega])}.
 \]
\end{lemma}

\begin{proof}
 See \texttt{cramer\_gelbrich.tex} for the full proof. The argument
 follows three steps: (i) Gelbrich formula for Gaussian measures on
 $L^2([0,T])$; (ii) operator square-root Lipschitz bound on operators
 bounded below by $\mu_{\min}$; (iii) Parseval's identity applied to
 the Cramér kernel $\Delta K(t,s) = \int (e^{i\omega t}-1)
 \overline{(e^{i\omega s}-1)} (S_1-S_2)(\omega)\,\mathrm{d}\omega$.
\end{proof}

\subsection{Hurst exponent consistency}

\begin{corollary}[Hurst recovery]\label{cor:hurst}
 If $\mathcal{L}(\theta) \leq \varepsilon$ and $u_{\mathrm{ref}}$
 has the power-law PSD $S_H(f) = C_H|f|^{1-2H}$, then the
 log-log OLS Hurst estimator $\hat{H}$ applied to $u_\theta$
 satisfies
 \[
  |\hat{H}(u_\theta) - H| \;\leq\; C_{\mathrm{OLS}}\sqrt{\varepsilon},
 \]
 where $C_{\mathrm{OLS}}$ depends on $[\omega_0, \Omega]$ but not on $\theta$.
\end{corollary}

\begin{proof}
 The proof is direct and does not invoke Theorem~\ref{thm:main}.

 \emph{Step 1 (Spectral control).}
 Since $S_H \leq M_{\max}$ on $[\omega_0,\Omega]$, the log-to-linear
 lemma \citep[Lemma~A.2]{TODO-paper1} gives
 $\|S_\theta - S_H\|_{L^2([\omega_0,\Omega])} \leq C_1\sqrt{\varepsilon}$.

 \emph{Step 2 (Variogram Parseval).}
 By the Cram\'er representation,
 $V(\tau) = \int|e^{i\omega\tau}-1|^2S(\omega)\,\mathrm{d}\omega$.
 Splitting at $\omega_0$ and applying Cauchy--Schwarz with
 $|e^{i\omega\tau}-1|^2\leq 4$:
 \[
  |V_\theta(\tau) - V_H(\tau)|
  \;\leq\; 4\sqrt{\Omega-\omega_0}\cdot C_1\sqrt{\varepsilon}
  + \tfrac{2C_H}{1-H}\,\omega_0^{2-2H}.
 \]

 \emph{Step 3 (OLS sensitivity).}
 The log-log OLS slope is Lipschitz in $\log V$: with weights
 $w_j = \log\tau_j - \bar{\ell}$ and
 $V_{\min} = c_H\tau_{\min}^{2H}$,
 \[
  |\hat H(u_\theta) - H|
  \;\leq\; \frac{\sum_j|w_j|}{V_{\min}\,\sum_j w_j^2}
   \cdot \max_j|V_\theta(\tau_j)-V_H(\tau_j)|.
 \]
 Substituting Step~2 and choosing
 $\omega_0 = O(\varepsilon^{1/(2-2H)})$ to absorb the tail term gives
 $|\hat H(u_\theta) - H| \leq C_{\mathrm{OLS}}\sqrt{\varepsilon}$
 with
 \[
  C_{\mathrm{OLS}} = \frac{4C_1\sqrt{\Omega-\omega_0}\,\sum_j|w_j|}
   {c_H\,\tau_{\min}^{2H}\,\sum_j w_j^2}.
 \]
 Full details and explicit constants in \texttt{cramer\_gelbrich.tex}.
\end{proof}

% ═════════════════════════════════════════════════════════════════════════════
\section{Experiments}
\label{sec:experiments}
% ═════════════════════════════════════════════════════════════════════════════

\subsection{Setup}

\paragraph{Model parameters.}
We use the rough Bergomi model \eqref{eq:rbergomi} with parameters
from \citet{gatheral2018volatility,bayer2016pricing}:
$H = 0.1$, $\eta = 1.9$, $\xi_0 = 0.04$ (initial vol 20\%),
$\rho = 0$ (leverage effect studied in \cref{sec:greeks}).
Time horizon $T = 1$ year, $N = 252$ daily steps.

\paragraph{Reference simulation.}
$10{,}000$ variance paths from Wood--Chan exact fBm simulation
\citep{wood1994simulation}, stock paths from Euler--Maruyama.

\paragraph{Surrogates.}
Three models trained for 3000 epochs on CPU (Adam, cosine LR schedule):
SpectralFilter (exact PSD), SpectralFilter (asymptotic PSD),
BochnerMLP (exact PSD).

\subsection{Hurst exponent recovery}

\begin{table}[h]
\centering
\caption{Hurst exponent recovery (structure function estimator),
Hurst sweep $H \in \{0.05, 0.10, 0.15, 0.20, 0.30, 0.50\}$.
SpectralFilter, 2000 epochs, smart initialisation.}
\label{tab:hurst}
\begin{tabular}{crrr}
\toprule
$H_{\mathrm{true}}$ & $\hat{H}_{\mathrm{exact}}$
 & $\hat{H}_{\mathrm{asymp}}$ & $\Delta H_{\mathrm{exact}}$ \\
\midrule
0.05 & 0.063 & 0.142 & $+0.013$ \\
0.10 & 0.109 & 0.171 & $+0.009$ \\
0.15 & 0.155 & 0.202 & $+0.005$ \\
0.20 & 0.201 & 0.237 & $+0.001$ \\
0.30 & 0.293 & 0.311 & $-0.007$ \\
0.50 & 0.473 & 0.471 & $-0.027$ \\
\bottomrule
\end{tabular}
\end{table}

Exact PSD training recovers $H$ with error $|\Delta H| \leq 0.013$
across all tested values. The asymptotic target systematically
over-estimates $H$ for $H \leq 0.3$ (bias up to $0.092$ for $H = 0.05$).
The $H = 0.5$ residual error ($-0.027$) is a finite-$N$ structural
artefact from the circulant periodicity and is the same for both targets.

\begin{figure}[t]
 \centering
 \includegraphics[width=0.62\textwidth]{hurst_recovery}
 \caption{%
  \textbf{Hurst exponent recovery.}
  Estimated $\hat{H}$ (structure-function OLS) vs.\ true $H$
  for the SpectralFilter trained on exact (blue) and asymptotic (orange)
  PSD targets. The diagonal (dashed) is perfect recovery.
  Exact PSD training lies within $|\Delta H| \leq 0.013$ of the diagonal
  across the full rough-vol range $H \in [0.05, 0.50]$.
  The asymptotic target systematically over-estimates $H$ for $H \leq 0.3$
  (bias up to $+0.09$ at $H=0.05$), consistent with Corollary~\ref{cor:hurst}
  and the mid-frequency PSD mis-match in Figure~\ref{fig:psd-comparison}.
 }
 \label{fig:hurst-recovery}
\end{figure}

\begin{remark}[Filter geometry: Bochner loss does not recover the Cholesky factor]
\label{rem:filter-geometry}
The SpectralFilter's learnable log-amplitudes $a_k$ do not converge to
the exact Cholesky factor $\tfrac{1}{2}\log(S_H(f_k)\cdot N/(2\,\mathrm{d}t))$
of the fGn covariance matrix.  Training filters at $H \in \{0.05, 0.10, 0.30\}$
on the exact fGn PSD and comparing the learned $a_k$ to the theoretical
Cholesky factor gives $\max_k|\Delta a_k| = 2.77$ in all cases
(Figure~\ref{fig:filter-geometry}).  The Bochner loss constrains the
\emph{expected} PSD of the output; many filter configurations satisfy
this in expectation without matching the per-frequency Cholesky amplitude.
Theorem~\ref{thm:main} is therefore a \emph{distributional} result:
the surrogate law converges to the reference law in $W_2$ but the filter
weights need not be the spectral Cholesky factor.
\end{remark}

\begin{figure}[h]
 \centering
 \includegraphics[width=\textwidth]{filter_geometry}
 \caption{%
  \textbf{Learned log-amplitude vs.\ theoretical Cholesky factor.}
  Blue: SpectralFilter learned $a_k$ after 3000 epochs.
  Orange: exact Cholesky factor $\tfrac{1}{2}\log(S_H(f_k)\cdot N/(2\,\mathrm{d}t))$.
  The maximum per-frequency gap $\max_k|\Delta a_k| = 2.77$ is identical
  for $H \in \{0.05, 0.10, 0.30\}$: the Bochner loss has a large solution
  manifold and the gradient descent finds a solution that satisfies the
  PSD constraint in expectation without converging to the Cholesky factor.
 }
 \label{fig:filter-geometry}
\end{figure}

\subsection{Hurst estimator comparison}
\label{sec:hurst-estimators}

We compare four estimators applied to $300$ SpectralFilter paths
($N=252$, $H \in \{0.05, 0.10, 0.20, 0.30\}$; same filter as
§\ref{sec:experiments}):
\begin{enumerate}[nosep,label=(\roman*)]
 \item \textbf{Structure function}: log-log OLS on
       $\mathbb{E}[(W_{t+\tau}-W_t)^2] \sim \tau^{2H}$; used throughout
       this paper.
 \item \textbf{PSD slope}: log-log OLS on the empirical PSD, fitting
       $\hat{S}(f) \sim f^{1-2H}$.
 \item \textbf{R/S analysis}: Hurst's rescaled-range,
       $\mathbb{E}[R(n)/S(n)] \sim n^H$.
 \item \textbf{Cram\'er exact}: grid-search
       $\hat{H}=\operatorname{arg\,min}_{H'}\|\log\hat{S}_\mathrm{emp}
       -\log S_{H'}\|^2$ over a pre-built grid
       $H' \in [0.02, 0.48]$ (step $0.01$, $10^4$ paths per node).
\end{enumerate}

\begin{table}[h]
\centering
\caption{Hurst estimator comparison: bias and standard deviation
($N=252$, 300~paths per $H$).
Structure function achieves smallest absolute bias ($|\text{bias}|\leq 0.011$
for $H\leq 0.30$).
Cram\'er exact achieves tightest standard deviation ($0.012$, nearly constant).
PSD slope and R/S are unreliable for rough vol: PSD slope incurs severe
negative bias at small $H$ ($-0.195$ at $H=0.05$); R/S has a large
positive bias ($+0.20$) uniform in $H$.}
\label{tab:hurst-estimators}
\begin{tabular}{lcccccccc}
\toprule
 & \multicolumn{2}{c}{$H=0.05$} & \multicolumn{2}{c}{$H=0.10$}
 & \multicolumn{2}{c}{$H=0.20$} & \multicolumn{2}{c}{$H=0.30$} \\
\cmidrule(lr){2-3}\cmidrule(lr){4-5}\cmidrule(lr){6-7}\cmidrule(lr){8-9}
Estimator & Bias & Std & Bias & Std & Bias & Std & Bias & Std \\
\midrule
Struct.\ function & $-0.002$ & $0.018$ & $-0.002$ & $0.028$ & $-0.006$ & $0.042$ & $-0.011$ & $0.052$ \\
PSD slope         & $-0.195$ & $0.159$ & $-0.120$ & $0.161$ & $-0.065$ & $0.162$ & $-0.007$ & $0.159$ \\
R/S               & $+0.208$ & $0.032$ & $+0.200$ & $0.037$ & $+0.173$ & $0.039$ & $+0.151$ & $0.047$ \\
Cram\'er exact    & $+0.055$ & $0.012$ & $+0.056$ & $0.012$ & $+0.055$ & $0.012$ & $+0.053$ & $0.011$ \\
\bottomrule
\end{tabular}
\end{table}

\begin{figure}[h]
 \centering
 \includegraphics[width=\textwidth]{hurst_estimator_comparison}
 \caption{%
  \textbf{Hurst estimator comparison.}
  Violin plots of $\hat{H}$ across 300 paths for each method.
  Structure function (blue) tracks the diagonal closely with small
  variance.  Cram\'er exact (green) is tightest but has a
  $+0.055$ positive bias from the finite-$N$ frequency cutoff.
  PSD slope (orange) collapses near zero for $H\leq 0.10$;
  R/S (red) overshoots by $+0.20$ uniformly.
 }
 \label{fig:hurst-estimators}
\end{figure}

The Cram\'er exact estimator has a systematic $+0.055$ positive bias
at all $H$, arising from the finite-$N$ frequency grid: the minimum
observable frequency $f_\mathrm{min} = 1/T = 1$~yr$^{-1}$ sits above
the low-frequency divergence of the fGn PSD, so the grid search finds
a slightly higher $H'$ to compensate for missing low-frequency energy.
A simple bias correction of $-0.055$ matches structure-function accuracy
while retaining the tighter standard deviation.

\textbf{Recommendation.} For rough vol ($H \in [0.05, 0.15]$) use the
structure function (smallest MAE).  PSD slope and R/S should \emph{not}
be used below $H < 0.2$: at $H=0.05$ both misestimate by multiples of
the true value.  The Cram\'er exact estimator provides the tightest
posterior (std~$=0.012$) and is the most robust to outlier paths, making
it preferable when a confidence interval on $H$ is required.

\subsection{Implied volatility surface}

\begin{table}[h]
\centering
\caption{IV surface RMSE (vol points) vs.\ MC reference.
Strikes $K/S_0 \in \{0.80, \ldots, 1.20\}$, maturities $T \in
\{0.25, 0.50, 1.00\}$. MC noise floor $\approx 1.0$~vp (10k paths).
Primary calibration: $H=0.1$, $\eta=1.9$, $\xi_0=0.04$.}
\label{tab:iv-rmse}
\begin{tabular}{lrrrr}
\toprule
Model & Overall & $T=0.25$ & $T=0.50$ & $T=1.00$ \\
\midrule
\multicolumn{5}{c}{\textit{$H=0.10$, $\eta=1.9$, $\xi_0=0.04$ (baseline)}} \\
SpectralFilter (exact PSD) & 1.33 & 1.16 & 0.66 & 1.89 \\
SpectralFilter (asymptotic) & 2.27 & 1.41 & 1.79 & 3.20 \\
BochnerMLP (exact PSD)   & \textbf{1.21} & 1.28 & 1.12 & \textbf{1.23} \\
\midrule
\multicolumn{5}{c}{\textit{$H=0.05$, $\eta=1.9$, $\xi_0=0.04$ (extreme roughness)}} \\
SpectralFilter (exact PSD) & \textbf{1.61} & \textbf{1.79} & \textbf{0.77} & \textbf{2.01} \\
SpectralFilter (asymptotic) & 2.89 & 1.94 & 2.29 & 4.01 \\
BochnerMLP (exact PSD)   & 2.03 & 2.13 & 1.82 & 2.13 \\
\midrule
\multicolumn{5}{c}{\textit{$H=0.10$, $\eta=1.5$, $\xi_0=0.04$ (mild vol-of-vol)}} \\
SpectralFilter (exact PSD) & 1.08 & 1.10 & 0.52 & 1.43 \\
SpectralFilter (asymptotic) & 1.64 & 1.00 & 1.26 & 2.34 \\
BochnerMLP (exact PSD)   & \textbf{0.61} & \textbf{0.60} & \textbf{0.53} & \textbf{0.69} \\
\midrule
MC noise floor ($10$k paths) & $\approx 1.0$ &  &  &  \\
\bottomrule
\end{tabular}
\end{table}

The exact PSD target yields a 1.7$\times$ lower overall RMSE than the
asymptotic approximation at the baseline parameters. The gap grows with
maturity: at $T=1$ the asymptotic model is $1.7\times$ worse (3.20 vs.\
1.89~vp). At extreme roughness ($H=0.05$) the gap widens further:
SpectralFilter~(asymptotic) reaches 4.01~vp at $T=1$ while the exact
target holds at 2.01~vp ($2\times$ better), confirming that the
power-law approximation is least accurate at mid-frequencies for small $H$.
Notably, at $H=0.05$ the SpectralFilter~(exact) outperforms the BochnerMLP
(1.61 vs.\ 2.03~vp overall): the nonlinear MLP struggles to match the
sharp roughness of extreme-$H$ paths, while the linear spectral filter is
optimal by construction. At lower vol-of-vol ($\eta=1.5$) the BochnerMLP
recovers strongly to 0.61~vp overall, below the 1.0~vp MC noise floor --
confirming that the spectral loss is architecture-agnostic and remains
effective across the calibrated parameter range.

\begin{figure}[t]
 \centering
 \includegraphics[width=\textwidth]{iv_surrogate_panel}
 \caption{%
  \textbf{Implied volatility smiles: MC reference vs.\ surrogates.}
  Each panel shows one maturity ($T \in \{0.25, 0.50, 1.00\}$);
  $x$-axis is moneyness $K/S_0$, $y$-axis is IV in \%.
  RMSE (vol points) relative to the MC reference is shown in the
  legend.
  SpectralFilter (exact PSD) tracks the reference closely at all
  maturities; SpectralFilter (asymptotic) incurs a systematic
  upward bias at $T=1.00$ (3.20~vp), where the mid-frequency
  PSD mis-match propagates to long-dated smile levels.
  BochnerMLP achieves the lowest per-maturity RMSE despite generating
  non-Gaussian paths, confirming that the spectral loss is
  architecture-agnostic.
 }
 \label{fig:iv-surrogate}
\end{figure}

\begin{figure}[t]
 \centering
 \includegraphics[width=\textwidth]{rmse_heatmap}
 \caption{%
  \textbf{IV RMSE (vol points) by strike and maturity.}
  Colour scale: green $=$ low error, red $=$ high error.
  Each cell shows RMSE relative to the 10k-path MC reference.
  SpectralFilter (exact PSD) and BochnerMLP maintain near-uniform
  low error across the entire strike--maturity grid.
  SpectralFilter (asymptotic) shows a concentrated error mass at
  long maturities ($T=1.00$) and deep OTM puts ($K/S_0 \leq 0.85$),
  precisely the region where the power-law PSD approximation is
  least accurate.
 }
 \label{fig:rmse-heatmap}
\end{figure}

\paragraph{PSD phase diagram.}
Figure~\ref{fig:psd-phase} quantifies the transition from exact-critical
to asymptotic-adequate regimes.  The relative $L^2(\log)$ PSD distance
between exact and asymptotic targets decreases monotonically with $H$
(from $0.072$ at $H=0.05$ to $0.011$ at $H=0.30$), consistent with
the power-law approximation becoming accurate as the PSD singularity
weakens.  The IV RMSE gap (exact vs.\ asymptotic) mirrors this:
$0.76$~vp gap at $H=0.05$, shrinking to $1.28$~vp gap at $H=0.30$
in the opposite direction (asymptotic actually deteriorates for large
$H$ where the filter overfits to the flat high-frequency region).
The critical threshold is $H \approx 0.15$, below which exact PSD
training is essential; above $H=0.20$ both targets yield comparable accuracy.

\begin{figure}[h]
 \centering
 \includegraphics[width=\textwidth]{psd_phase_diagram}
 \caption{%
  \textbf{PSD phase diagram.}
  Left: relative $L^2(\log)$ distance between exact and asymptotic fGn PSD
  as a function of $H$.  Right: IV RMSE vs.\ MC reference for SpectralFilter
  trained on exact (blue) and asymptotic (orange) PSD.
  Below $H \approx 0.15$ (shaded) the exact target is essential;
  above $H = 0.20$ both targets achieve comparable accuracy.
  Baseline parameters: $\eta=1.9$, $\xi_0=0.04$, $\rho=0$, $N=252$.
 }
 \label{fig:psd-phase}
\end{figure}

\subsection{Observation-horizon sensitivity}
\label{sec:horizon}

The minimum observable frequency in the training PSD is
$f_\mathrm{min} = 1/T$, so longer horizons resolve more of the
low-frequency region near the fGn PSD singularity.  We sweep
$T \in \{0.5, 1, 2, 5\}$~years at fixed $H=0.1$,
$\Delta t = 1/252$, using $N = \lfloor T\times 252\rfloor$ steps
and evaluating IV RMSE at terminal maturity $T$.

\begin{table}[h]
\centering
\caption{Observation-horizon sensitivity ($H=0.1$, $\Delta t=1/252$,
3 evaluation seeds).  IV RMSE at terminal maturity $T$;
$|\Delta H|$ from structure-function estimator.
Surrogate performs well for $T\leq 2$~yr; IV RMSE spikes at
$T=5$~yr due to the amplified low-frequency cutoff penalty
$C_3\omega_0^{-(1-2H)/2}$ (Theorem~\ref{thm:main}) and the
reduced MC budget at large $N$.}
\label{tab:horizon-sweep}
\begin{tabular}{lrrrc}
\toprule
$T$ (yr) & $N$ & $f_\mathrm{min}$ (Hz) & $|\Delta H|$ & IV RMSE (vp) \\
\midrule
0.5 & 126  & 2.000 & $0.010$ & $1.36\pm 0.34$ \\
1.0 & 252  & 1.000 & $0.010$ & $1.57\pm 0.39$ \\
2.0 & 504  & 0.500 & $0.006$ & $1.32\pm 0.75$ \\
5.0 & 1260 & 0.200 & $0.003$ & $3.72\pm 1.03$ \\
\bottomrule
\end{tabular}
\end{table}

\begin{figure}[h]
 \centering
 \includegraphics[width=\textwidth]{horizon_sweep}
 \caption{%
  \textbf{Horizon sensitivity.}
  Left: mean absolute Hurst estimation error $|\hat H - H|$ vs.\ horizon $T$.
  Right: IV RMSE at terminal maturity vs.\ $T$.
  Hurst error decreases monotonically with $T$ (more data $\Rightarrow$ better
  frequency resolution).  IV RMSE is non-monotone: it stays low for
  $T\leq 2$~yr and spikes at $T=5$~yr, where the low-frequency singularity
  degrades the PSD fit.
 }
 \label{fig:horizon-sweep}
\end{figure}

Two effects compete as $T$ grows.  Longer horizons reduce $|\Delta H|$
monotonically ($0.010 \to 0.003$), consistent with Corollary~\ref{cor:hurst}.
But they also move $f_\mathrm{min}$ closer to the PSD singularity at $f=0$,
amplifying the cutoff term $C_3\omega_0^{-(1-2H)/2}$ in
Theorem~\ref{thm:main}.  At $T=5$~yr the IV RMSE spikes to
$3.72\pm 1.03$~vp, with large cross-seed variance signalling instability
in the low-frequency PSD fit.  We therefore recommend the Cramér PINN for
$T\leq 2$~yr, where IV RMSE stays below $1.6$~vp at the primary calibration
($H=0.1$, $\eta=1.9$).

\subsection{Runtime comparison}

\begin{table}[h]
\centering
\caption{Wall-clock time for generating $10{,}000$ variance paths
($N=252$, $T=1$ year). Surrogates benchmarked on NVIDIA GTX~950 (2~GB);
Wood--Chan runs on CPU (single core).
Training cost amortised over $10^6$ generated paths (${}^\dagger$estimated).}
\label{tab:timing}
\begin{tabular}{lrrr}
\toprule
Method & Train (s) & Gen $10$k (ms) & Amort.\ $10^6$ (ms/10k) \\
\midrule
Wood--Chan (exact fBm)   &  & 365  & 365   \\
Bennedsen hybrid scheme   &  & ${}^\dagger$200 & ${}^\dagger$200 \\
SpectralFilter (exact PSD) & 5.7 & 0.26 & 57   \\
SpectralFilter (asymptotic) & 5.7 & 0.26 & 57   \\
BochnerMLP (exact PSD)   & 6.5 & 0.17 & 66   \\
\midrule
3-factor Markovian Bergomi &  & ${}^\dagger$30 & ${}^\dagger$30 \\
Lifted Heston ($n=5$ factors)&  & ${}^\dagger$10 & ${}^\dagger$10 \\
\bottomrule
\end{tabular}
\end{table}

Training is a one-time cost; amortised over $10^6$ paths the spectral
surrogates generate each batch of $10{,}000$ paths in under 1\% of the
Wood--Chan time. At $10^7$ paths (deep hedging scale), the amortised
cost falls to $\approx 6$~ms/10k for SpectralFilter, a $60\times$
wall-clock speedup over exact fBm simulation. SpectralFilter generates
paths via a single RFFT/IRFFT round-trip; BochnerMLP uses a forward pass
through a 35k-parameter MLP, which on GPU is slightly faster than the FFT
chain due to fused matrix-multiply kernels.

\paragraph{Reference IV surface ($H=0.1$, $\rho=0$).}
\begin{center}
\begin{tabular}{lrrrrrrrrr}
\toprule
$T \backslash K/S_0$ & 0.80 & 0.85 & 0.90 & 0.95 & 1.00 & 1.05 & 1.10 & 1.15 & 1.20 \\
\midrule
0.25 & 25.9 & 22.8 & 20.1 & 18.1 & 17.2 & 18.0 & 19.6 & 21.4 & 23.3 \\
0.50 & 23.2 & 20.8 & 18.8 & 17.3 & 16.8 & 17.4 & 18.6 & 19.9 & 21.2 \\
1.00 & 22.6 & 20.6 & 19.0 & 17.9 & 17.4 & 17.7 & 18.5 & 19.5 & 20.5 \\
\bottomrule
\end{tabular}
\end{center}
IV in \%. ATM vol $\approx 17\%$ (slightly below $\sqrt{\xi_0} = 20\%$
due to Euler--Maruyama discretisation bias). Smile is bowl-shaped and
symmetric (expected for $\rho=0$); term structure nearly flat.

\subsection{CausalSpectralFilter: supporting $\rho \neq 0$}
\label{sec:causal}
% ═════════════════════════════════════════════════════════════════════════════

The experiments in \cref{sec:greeks} establish that the irreducible
$-0.27$ delta bias originates from missing the joint $(V,S)$ law. A
natural prior question is whether the spectral surrogate can at least
\emph{encode} $\rho \neq 0$ correctly within rBergomi, so that the
leverage channel is present even if the exact $S$--$V$ law is not matched.

\paragraph{Setup: $z$-Cholesky and the BLP reference.}
The \emph{$z$-Cholesky} approach reuses the filter's i.i.d.\ input noise
$z \sim \mathcal{N}(0,I)$ as the stock BM driver:
\begin{equation}
 d\!\log S_k
  = -\tfrac{1}{2}V_k\,dt
   + \sqrt{V_k\,dt}\,\bigl(\rho\,z_k + \sqrt{1-\rho^2}\,z_k^\perp\bigr).
 \label{eq:z-cholesky}
\end{equation}
The ground truth is an RL-discretised rBergomi path with $\rho = -0.9$,
simulated via the \emph{BLP hybrid scheme} \citep{bennedsen2017hybrid}
(same Volterra kernel as \cref{sec:greeks}, now applied to rBergomi).
The BLP reference produces a left-skewed IV surface (ATM $\approx 15\%$
at $T=0.25$, declining to $9\%$ at $K/S_0 = 1.10$) and ATM delta
$0.67$--$0.71$, consistent with the leverage effect.

\paragraph{Circular filter: catastrophic failure.}
Applying \eqref{eq:z-cholesky} to the circular SpectralFilter gives an
IV RMSE of $204$~vp and ATM deltas rising from $1.07$ to $1.90$ --
physically impossible for European calls. The mechanism is
\emph{non-causality of the circular FFT}. The filter computes
\[
 \Delta\hat{W}_k
  = \bigl[\mathrm{IRFFT}\!\bigl(\tilde{H}\cdot\mathrm{RFFT}(z)\bigr)\bigr]_k,
\]
a circular convolution over the \emph{entire} horizon $N$. Consequently
$V_k = \xi_0 e^{\eta\hat{W}_k - \cdots}$ depends on future innovations
$z_{k+1},\ldots,z_{N-1}$ through the circular wrap-around, so $V_k$ and
$d\!\log S_{k+j}$ are spuriously correlated forward in time, violating
the martingale condition.

\paragraph{CausalSpectralFilter: architecture and initialisation.}
We replace the circular FFT with a causal lower-triangular Toeplitz
convolution:
\begin{equation}
 \Delta\hat{W}_k = \sum_{j=0}^{k} h_{k-j}\,z_j,
 \label{eq:causal-filter}
\end{equation}
where $h \in \mathbb{R}^N$ is a learnable real impulse response.
Because output at step $k$ depends only on $z_0,\ldots,z_k$, we have
$V_k \perp z_k$ and \eqref{eq:z-cholesky} is mathematically valid.
We implement \eqref{eq:causal-filter} as a single \texttt{F.conv1d}
with $(N-1)$ left-zero-padding, costing one matrix–vector product per
forward pass.

Initialisation uses the Riemann--Liouville increment kernel:
\begin{equation}
 h_0 = \frac{dt^H}{\Gamma(H+\tfrac{1}{2})},
 \qquad
 h_m = \frac{dt^H\bigl(g(m+1)-g(m)\bigr)}{\Gamma(H+\tfrac{1}{2})},
 \quad m \geq 1,
 \label{eq:rl-init}
\end{equation}
where $g(i)=i^{H-1/2}$. For $H < \tfrac{1}{2}$, $h_m < 0$ for $m \geq 1$
(roughness is encoded as negative memory in the innovations). We then
fine-tune with the same Bochner PSD loss used throughout, reaching a
final log-PSD MSE of $0.019$ (vs.\ $0.017$ for the circular filter) in
$3\,000$ epochs. The CausalSpectralFilter has $N = 252$ parameters
(vs.\ $N/2 = 126$ for the circular architecture) and an inference cost
of $12.6$~ms per $10\,000$ paths (vs.\ $1.9$~ms for the circular FFT)
on the same GTX 950 GPU, a $6.7\times$ slowdown that remains
$29\times$ faster than the Wood--Chan reference simulator. Hurst
recovery via structure function is preserved
($|\hat{H}-H| = 0.003$ for causal vs.\ $0.009$ for circular, both well
within the $\sqrt{\varepsilon}$ bound of \cref{cor:hurst}).

\paragraph{Results.}
Table~\ref{tab:causal-blp} and Figure~\ref{fig:causal-iv} summarise.
The causal filter with $z$-Cholesky achieves an IV RMSE of
\textbf{3.26~vp}, a \textbf{63$\times$ improvement} over the circular
architecture and better than the $\rho=0$ causal baseline (6.8~vp).
All ATM deltas are in $[0,1]$; the ATM bias shrinks from $-0.17$
($\rho=0$ baseline) to $-0.07$ ($z$-Cholesky), recovering $\sim\!60\%$
of the leverage-effect contribution.

\begin{table}[h]
\centering
\caption{IV RMSE vs BLP rBergomi reference ($\rho=-0.9$, 8k paths each)
and ATM $\Delta$ at $T=0.50$. Causal $z$-Cholesky is 63$\times$ better
than the circular architecture and reduces the delta bias by 60\% vs.\
the $\rho=0$ baseline.}
\label{tab:causal-blp}
\begin{tabular}{lrrrr}
\toprule
Condition & \multicolumn{3}{c}{IV RMSE (vp) per maturity} & ATM $\Delta$ \\
     & $T=0.25$ & $T=0.50$ & $T=1.00$ & $T=0.50$ \\
\midrule
BLP reference ($\rho=-0.9$)      &   &   &   & 0.689 \\
Circular SF  z-Cholesky ($\rho=-0.9$) & 57.7 & 89.8 & 359.6 & $>1$ \\
Causal SF   baseline  ($\rho=0$)  & 7.6 & 7.5 &  5.3 & 0.517 \\
Causal SF   z-Cholesky ($\rho=-0.9$) & \textbf{4.1} & \textbf{3.7} & \textbf{1.4} & \textbf{0.620} \\
\bottomrule
\end{tabular}
\end{table}

\begin{figure}[t]
 \centering
 \includegraphics[width=\textwidth]{causal_filter_comparison}
 \caption{%
  \textbf{CausalSpectralFilter vs circular SpectralFilter ($\rho=-0.9$).}
  Each panel shows one maturity.
  BLP reference (black) has the correct left skew.
  Causal $z$-Cholesky (green dashed) tracks the reference closely at
  all strikes; causal $\rho=0$ baseline (blue dotted) matches the IV
  \emph{level} but not the skew. Circular $z$-Cholesky (red
  dash-dot) is physically invalid (IVs exceed 100\% and deltas $>1$
  at long maturities, not shown on this scale).
 }
 \label{fig:causal-iv}
\end{figure}

\paragraph{Robustness across $H$.}
The improvement is not specific to $H = 0.1$.
Table~\ref{tab:causal-hurst} reports the sweep across the typical
rough-vol range $H \in \{0.05, 0.10, 0.15\}$. The causal $z$-Cholesky
IV RMSE stays in $[2.4, 3.5]$~vp across all $H$, while the circular
baseline exceeds $100$~vp at every $H$. The improvement factor
varies between $29\times$ and $60\times$, peaking at $H=0.1$ where the
circular failure is most severe. ATM delta bias decreases monotonically
with $H$ (from $-0.076$ at $H=0.05$ to $-0.050$ at $H=0.15$), consistent
with smoother $V$ paths being easier to match in distribution.

\begin{table}[h]
\centering
\caption{CausalSpectralFilter robustness sweep across $H$, all conditions
at $\rho=-0.9$, 8k paths. IV RMSE in vol points vs BLP reference;
$\Delta$ bias = $\hat\Delta_{\mathrm{ATM}} - \Delta_{\mathrm{ATM,ref}}$ at $T=0.50$.}
\label{tab:causal-hurst}
\begin{tabular}{ccccccc}
\toprule
$H$ & Circular SF & Causal SF & Causal SF & Improvement & ATM $\Delta$ bias \\
  & z-Cholesky & $\rho=0$ & z-Cholesky & factor   & (causal $z$-Chol) \\
\midrule
0.05 & 103.0 & 7.10 & 3.54 & $29\times$ & $-0.076$ \\
0.10 & 204.6 & 6.84 & \textbf{3.41} & $\mathbf{60\times}$ & $-0.073$ \\
0.15 & 106.1 & 6.83 & 2.39 & $44\times$ & $-0.050$ \\
\bottomrule
\end{tabular}
\end{table}

\paragraph{Closing the residual: matched-PSD training.}
The remaining $\sim\!-0.07$ delta bias of the causal $z$-Cholesky
surrogate is attributable to the V-distribution gap between the
RL-discretised BLP reference and the Wood--Chan-based fGn PSD used in
training. We test this directly by training a second causal filter on
the \emph{RL-fBm} PSD (estimated by Monte Carlo from $10^5$ paths of the
same RL discretisation that drives the reference). Table~\ref{tab:rl-match}
shows the result:

\begin{table}[h]
\centering
\caption{Effect of matching the training PSD to the reference V
distribution ($H=0.1$, $\rho=-0.9$, 5k surrogate paths). IV RMSE is
against the BLP reference; bias is
$\hat\Delta_{\mathrm{ATM}} - \Delta_{\mathrm{ATM,ref}}$. Matching the
training PSD to the RL-fBm distribution eliminates the residual bias.}
\label{tab:rl-match}
\begin{tabular}{lcccc}
\toprule
Training PSD & IV RMSE (vp) & Bias $T=0.25$ & Bias $T=0.50$ & Bias $T=1.00$ \\
\midrule
Wood--Chan (exact fBm)     & 3.42 & $-0.086$ & $-0.073$ & $-0.058$ \\
RL-fBm (matched to reference) & \textbf{2.57} & \textbf{$+0.011$} & \textbf{$+0.011$} & \textbf{$-0.007$} \\
\bottomrule
\end{tabular}
\end{table}

The matched-PSD surrogate achieves $2.57$~vp and an ATM delta bias of
$\lesssim 0.011$ at all maturities, essentially Monte Carlo noise.
This confirms the diagnosis: the leverage-channel residual was an
artefact of the training-vs-reference V-distribution mismatch, not a
deficiency of the causal architecture. The lesson generalises: for
any reference simulator, training the causal SpectralFilter on the
simulator's empirical PSD recovers $\rho \neq 0$ Greeks faithfully.

\paragraph{Two negative results: spectral matching of rough Heston $V$
makes things worse.}
We attempted two natural variants of the causal SpectralFilter that
match rough Heston's $V$ marginal more closely:
\emph{(i)} train on the PSD of $V_{\mathrm{rH}}-\theta$ and use a linear
reconstruction $V_{\mathrm{surr}} = \theta + \mathrm{filter}(z)$;
\emph{(ii)} train on the PSD of $\log V_{\mathrm{rH}}$ and use the
rBergomi-style exp reconstruction
$V_{\mathrm{surr}} = \theta \exp(W - \tfrac{1}{2}\mathrm{Var}\,W)$.
Both \textbf{worsen} the cross-model bias (Table~\ref{tab:rh-negatives}).
Variant~(i) requires clipping $44\%$ of cells to enforce positivity,
inflating $\E[V_{\mathrm{surr}}]$ to $0.15$ ($3.7\times$ target).
Variant~(ii) collapses $\E[V_{\mathrm{surr}}]\to 0.004$ because the
empirical variance correction is dominated by
$\mathrm{Var}(d\log V_{\mathrm{rH}}) \approx 62$, log-$V$ excursions
when $V_{\mathrm{rH}}$ briefly visits very small values.

\begin{table}[h]
\centering
\caption{Attempts to match rough Heston $V$ within the spectral framework
all underperform the rBergomi-$V$ baseline. Spectral surrogates capture
second-order statistics; rough Heston $V$'s skewness and heavy tails
require higher-order modelling.}
\label{tab:rh-negatives}
\begin{tabular}{lccc}
\toprule
Surrogate $V$ design & IV RMSE (vp) & ATM $\Delta(T=.5)$ & Bias \\
\midrule
Causal SF, rBergomi $V$ (Surr~D above)       & \textbf{4.57} & \textbf{0.6247} & \textbf{$-0.150$} \\
Causal SF, linear $V = \theta + \mathrm{filter}(z)$ & 20.4 & 0.456 & $-0.319$ \\
Causal SF, $\log V$ matched, exp recon       & 9.7 & 0.485 & $-0.290$ \\
\bottomrule
\end{tabular}
\end{table}

Both failures reinforce the diagnosis: rough Heston's $V$ marginal is
too skewed and heavy-tailed for Gaussian second-order matching, and the
rBergomi log-normal $V$, \emph{despite a wrong distribution}, is the
best base process within the spectral family. Closing the remaining
$-0.15$ bias requires a Volterra-aware (non-spectral) architecture
that captures rough Heston's mean reversion and tail structure.

\paragraph{Sub-quadratic Volterra updates.}
The Volterra recursion above is $O(N^2)$ per path because the kernel
sum at step $k$ uses all of $V_0, \ldots, V_k$. Approximating the RL
kernel by a sum of $M$ exponentials,
$w_i \approx \sum_{m=1}^{M} a_m\, r_m^{\,i}$ with $r_m = e^{-\lambda_m}$,
turns the sum into $M$ recursively updatable states:
\begin{equation}
 S_m^X(k) = X_k + r_m\, S_m^X(k-1),
 \qquad \text{cost: } O(M) \text{ per step.}
 \label{eq:markov-state}
\end{equation}
We fit $(a_m, \lambda_m)$ from the rough-Heston RL kernel by
geometrically spacing the $\lambda_m$ over $[1/N, 1]$ and solving for
$a_m$ by nonneg least squares. The resulting
\emph{MarkovianVolterraSurrogate} has the same V-conditional structure
\eqref{eq:volterra-surrogate} but $O(MN)$ per path.
Table~\ref{tab:markov-volterra} shows the sweep:

\begin{table}[h]
\centering
\caption{Sub-quadratic Volterra via $M$-mode RL kernel approximation
($H=0.1$, $N=252$, batch $5\,000$, GPU). $M = 5$ already saturates the
kernel-fit error ($\approx 4\%$ relative) and gives the best IV-RMSE
of the sweep, $7\times$ faster than the exact $O(N^2)$ baseline.}
\label{tab:markov-volterra}
\begin{tabular}{lcccc}
\toprule
Surrogate & Fwd time (ms) & Speedup & IV RMSE (vp) & ATM $\Delta$ bias \\
\midrule
$O(N^2)$ Volterra (exact)      & 280 & $1\times$  & 0.91 & $-0.004$ \\
Markovian, $M=3$           & 47 & $6.0\times$ & 0.76 & $-0.001$ \\
\textbf{Markovian, $M=5$}      & \textbf{39} & \textbf{$7.2\times$} & \textbf{0.73} & \textbf{$+0.001$} \\
Markovian, $M=7$           & 40 & $7.1\times$ & 0.88 & $-0.008$ \\
Markovian, $M=10$          & 40 & $7.1\times$ & 0.85 & $-0.009$ \\
\bottomrule
\end{tabular}
\end{table}

All Markovian variants match the exact baseline within MC noise; the
$7\times$ speedup is bounded below by Python loop overhead and would
grow as $O(N/M)$ for longer time horizons. The Greeks are robust to a
$4\%$ relative kernel-fit error, a strong empirical indicator that
sub-quadratic approximations are practically usable.

\paragraph{Non-affine generalisation and full joint training.}
The diffusion term $\sqrt{V_j}$ in \eqref{eq:volterra-surrogate} is a
Heston-specific choice ($\sigma(V) = \nu V^{1/2}$). The same
autoregressive form accepts $\sigma(V) = V^\beta$ with $\beta$
learnable, call this the \emph{PowerVolterraSurrogate}. For
$\beta = 1$ this is a rough-lognormal-vol process; for $\beta = 0$ a
rough additive-vol process; intermediate $\beta$ interpolates.
Training the full parameter set
$(\theta, c_{\rm drift}, c_{\rm stoch}, \beta, \{w_i\})$ with the
price-plus-moment loss from a heavily misspecified initial state
$(\theta, \kappa, \nu, \beta) = (0.10, 1.0, 0.5, 0.8)$ \emph{and} a
$20\%$-perturbed kernel reaches IV RMSE $1.85$~vp and ATM delta bias
$+0.014$, closer to the RH-initialised $0.80$~vp ceiling than to
the spectral-floor $4.57$~vp. Final learned $\beta = 0.60$ (target
$0.5$). Individual parameter values do not match rough Heston, but
the joint trajectory and observables do, classical price-only
identifiability that the moment loss only partially resolves.

\paragraph{Live SPY option chain calibration (multi-maturity).}
We calibrate the joint Volterra surrogate to real SPY mid prices pulled
from Yahoo Finance at three expirations simultaneously: $T = 32$, $88$,
and $179$ days, with $S_0 = \$745.64$. All parameters
$(\theta, c_{\mathrm{drift}}, c_{\mathrm{stoch}}, \beta, \{w_i\})$
are \emph{shared} across the three maturities, only one $V$ path is
generated, and prices at each $T_i$ read from the appropriate column of
$V$. Twenty-six liquid OTM/ATM/ITM strikes total span
$K/S_0 \in [0.85, 1.15]$; market bid/ask spread implies $\sim\!1.4$~vp
IV uncertainty. Initialised at
$\theta_0 = (\mathrm{IV}^{\mathrm{ATM, short}})^2$ with $\rho = -0.85$,
the surrogate converges in $500$ Adam steps ($3$~min, GTX~950) to:

\begin{table}[h]
\centering
\caption{Multi-maturity calibration to live SPY options
($S_0 = \$745.64$, $\rho = -0.85$, $26$ calibration points across three
expirations). Overall IV RMSE $3.97$~vp; the short-end fits within
$2\times$ bid/ask uncertainty, while the long-end shows a systematic
under-pricing of OTM calls characteristic of single-kernel rough-vol
calibrations (the surrogate's variance grows too slowly with $T$).}
\label{tab:spy-multi}
\begin{tabular}{lccc}
\toprule
Expiry & $T$ (yr) & Strikes & IV RMSE (vp) \\
\midrule
2026-06-26 ($+32$d) & 0.088 & 8 & \textbf{2.27} \\
2026-08-21 ($+88$d) & 0.241 & 9 & 3.79 \\
2026-11-20 ($+179$d) & 0.490 & 9 & 5.17 \\
\midrule
\multicolumn{3}{l}{\textbf{Overall}} & \textbf{3.97} \\
\bottomrule
\end{tabular}
\end{table}

Final learned parameters
$(\theta, c_{\mathrm{stoch}}, \beta) = (0.037, 0.72, 0.52)$ recover
near-Heston dynamics ($\beta \approx 0.5$); the $\theta$ value is
consistent with the long-run SPX VIX level. The term-structure
residual at the $179$-day expiry, a $\sim 5$~vp uniform downward bias
across all strikes, reflects a known limitation of \emph{stationary-
parameter} rough-vol calibrations: a single kernel cannot match the
SPX ATM IV term structure (rising from $15.4\%$ at $32$d to $20.8\%$ at
$179$d here) without time-varying forward variance. Closing this is
the multi-maturity forward-consistency direction listed below.

\paragraph{SPX Hurst index from ATM term structure.}
Before the structural calibration, we estimate the implied Hurst exponent
directly from the SPX ATM IV term structure at three maturities
($T \in \{0.088, 0.241, 0.490\}$~yr, ATM IVs $15.6\%$, $18.1\%$, $20.0\%$).
Under rBergomi, $\mathrm{ATM\text{-}IV}(T) \propto T^{H-1/2}$ (Fukasawa
approximation), so a log-log OLS fit gives
$\hat{H}_\mathrm{SPX} = 0.145 \pm 0.002$ (standard error across 500
bootstrap resamples of the three-point fit).
We then train a SpectralFilter at $H = 0.145$ and evaluate against a BLP
reference at the same parameters: $H_\mathrm{est} = 0.151$
(structure-function, $|\Delta H| = 0.007$), IV RMSE $= 1.43$~vp
(Figure~\ref{fig:spx-hurst}), and variance ratio
$\mathrm{Var}(V_\mathrm{surr})/\mathrm{Var}(V_\mathrm{ref}) = 0.51$
(some variance is absorbed in the Itô correction).
The $|\Delta H| = 0.007$ is within the $\sqrt{\varepsilon}$ bound of
Corollary~\ref{cor:hurst} and confirms that the Cramér PINN can be
fitted directly to the market-implied Hurst index with no additional
calibration.

\begin{figure}[h]
 \centering
 \includegraphics[width=\textwidth]{spx_hurst_analysis}
 \caption{%
  \textbf{SPX Hurst analysis.}
  Left: log-log ATM IV vs.\ maturity; slope gives
  $\hat{H}_\mathrm{SPX} = 0.145$.
  Right: IV surface, SpectralFilter (trained at $H=0.145$, $\eta=1.485$)
  vs.\ BLP reference ($T\in\{0.25,0.50,1.00\}$~yr,
  $K/S_0\in[0.80,1.20]$).  IV RMSE $= 1.43$~vp.
 }
 \label{fig:spx-hurst}
\end{figure}

\paragraph{Forward-variance parameterisation: rBergomi-style.}
Multi-maturity SPY calibration with stationary $\theta$ (above) leaves
a uniform $\sim\!5$~vp residual at the long expiry, driven by SPX's
rising ATM-IV term structure ($15.4\% \to 20.8\%$ from $32$d to $179$d)
that a scalar $\theta$ cannot match. The cleanest structural fix is
the rBergomi forward-variance parameterisation:
\begin{equation}
 V_t \;=\; \xi(t) \cdot \exp\!\Bigl(\eta\, \hat W^H_t
     \;-\; \tfrac{1}{2}\eta^2\,\mathrm{Var}(\hat W^H_t)\Bigr),
 \label{eq:fwd-var}
\end{equation}
where $\hat W^H_t = \mathrm{cumsum}(\mathrm{CausalSF}(z))_t$ and $\xi(t)$
is a piecewise-linear forward-variance curve parameterised by knots at
the calibration maturities (log-parameterisation enforces positivity).
By the lognormal martingale identity, $\mathbb{E}[V_t] = \xi(t)$
exactly, so the ATM-IV term structure is pinned to $\sqrt{\xi(t)}$ by
construction. V is positive without clamping, the Feller-violation
pathology that broke the time-varying $\theta(t)$ Volterra attempt
vanishes.

Calibrating on the same SPY chain (with an empirical-mean renormalisation
that enforces $\mathbb{E}[V_t] = \xi(t)$ exactly at the finite batch),
we reach IV RMSE \textbf{$5.25$~vp} overall. The per-maturity breakdown
clarifies the residual:
\begin{table}[h]
\centering
\caption{rBergomi forward-variance calibration on SPY (same chain as
Table~\ref{tab:spy-multi}, $\rho=-0.85$, $\eta = 1.485$). Martingale
property holds exactly after the empirical-mean correction; remaining
residual is moneyness-driven, not maturity-driven.}
\label{tab:spy-rbergomi}
\begin{tabular}{lcccc}
\toprule
Maturity & ATM IV err & Deep-ITM err ($K/S_0 = 0.85$) & Deep-OTM err & Per-mat RMSE \\
\midrule
32d & $-1.2$~vp & $-13.7$~vp & $-5.1$~vp & 6.13 vp \\
88d & $-2.0$~vp & $-8.5$~vp & $-1.4$~vp & 4.05 vp \\
179d & $-4.8$~vp & $-8.7$~vp & $-2.1$~vp & 5.46 vp \\
\bottomrule
\end{tabular}
\end{table}

ATM strikes fit to within $\sim\!2$~vp; the residual is dominated by
\emph{deep ITM} ($K/S_0 \leq 0.90$) where the surrogate's smile is too
flat, \textbf{this is a skew mismatch, not a term-structure or
positivity issue.} A $\$2$ price error at $K=\$635$ (where the option
is mostly intrinsic value) is $10\%$ of the time value and only $2\%$
of the total price, but corresponds to $\sim\!14$~vp of IV error
because deep-ITM vega is tiny. Adam-with-uniform-price-MSE cannot see
this in its gradients.

\paragraph{Remaining open directions for Paper~3:}
The structural calibration story is now closed: stationary-$\theta$
Volterra ($3.97$~vp, no term structure), time-varying-$\theta$
Volterra ($5.59$~vp, Feller violation), and rBergomi forward variance
($5.25$~vp, no Feller, exact martingale). Closing the remaining
moneyness-driven residual requires:
(i) IV-direct (vega-weighted or implicit-differentiation) loss to make
the optimiser see deep-strike error;
(ii) two-stage calibration (ATM curve first, then smile) so $\eta$ can
be pushed to SPX-realistic levels without breaking ATM fit;
(iii) richer auxiliary moment losses (autocovariance, third moments) to
constrain the joint $(V, S)$ law beyond second-order matching.

% ═════════════════════════════════════════════════════════════════════════════
\section{Discussion}
\label{sec:discussion}
% ═════════════════════════════════════════════════════════════════════════════

\paragraph{Why the exact PSD matters.}
The power-law PSD $S_H(f) \sim C|f|^{1-2H}$ is the correct
\emph{asymptotic} behaviour of fGn for $f \to 0$. At mid-frequencies
the exact PSD deviates from the power law; for $H = 0.1$ this deviation
exceeds 10\% relative error across most of the frequency range.
Matching the asymptotic shape correctly recovers the power-law slope but
mis-trains the correlation structure at short lags, biasing long-maturity
option prices (the $T = 1$ column in Table~\ref{tab:iv-rmse}).
The PSD phase diagram (Figure~\ref{fig:psd-phase}) quantifies the
regime boundary: exact PSD training is essential below $H\approx 0.15$;
above $H=0.20$ both targets converge to comparable IV RMSE. The
relative $L^2(\log)$ PSD distance drops from $0.072$ at $H=0.05$
to $0.011$ at $H=0.30$, matching the observed IV RMSE gap.
Horizon sensitivity (Table~\ref{tab:horizon-sweep}) shows this regime
boundary also applies in the time domain: longer observation windows
help Hurst recovery but amplify the low-frequency cutoff penalty,
making the Cramér PINN most reliable for $T \leq 2$~yr.

\paragraph{BochnerMLP: architecture-agnostic spectral loss.}
The MLP surrogate outperforms SpectralFilter on overall RMSE despite
generating non-Gaussian paths. This is expected: the MLP has more
capacity (41k vs.\ 126 parameters) and can approximate the fGn PSD more
closely. However, non-Gaussianity means the Itô correction
$-\frac{1}{2}\eta^2 t^{2H}$ does not exactly ensure $\E[\hat{V}_t] =
\xi_0$ (Jensen's inequality with non-Gaussian $\hat{W}_t$); we use the
empirical variance correction \eqref{eq:vhat} to mitigate this.
SpectralFilter, being a linear filter of Gaussian input, is
\emph{guaranteed} to produce Gaussian fGn paths and avoids this issue
entirely.

\paragraph{Application to deep hedging.}
\citet{buhler2019deep} require large batches of simulated price paths for
training a hedging neural network. The spectral surrogate reduces
path generation from $O(N^2)$ (Cholesky-based fBm) or $O(N \log N)$
(Wood--Chan) to a single forward pass of the trained network, enabling
on-the-fly batch sampling. The $\approx 1.2$~vp IV accuracy is well
within the calibration precision of typical options desks.
However, \cref{sec:greeks} shows that Greeks require the joint $(V, S)$
law: the spectral surrogate, which matches only the marginal distribution
of $V$, carries an irreducible delta bias of $\approx -0.27$ at ATM
when $\rho = -0.9$. Deep hedging with a spectral surrogate therefore
requires either a joint-law surrogate (future work) or a Malliavin
score correction to recover the leverage contribution. A circular-FFT $z$-Cholesky attempt fails (204~vp, delta $>1$) due to
wrap-around making $V_k$ depend on future innovations
(Section~\ref{sec:greeks}). The CausalSpectralFilter resolves this
exactly, achieving 3.26~vp and delta bias $-0.07$, a 63$\times$ RMSE
improvement available at no additional training cost.

\paragraph{Comparison to Markovian approximations.}
A standard industry approach to rough volatility is to approximate the
fractional kernel by a finite sum of exponentials, yielding a
\emph{Markovian} multi-factor model that can be simulated with
conventional Euler schemes. Two leading constructions are:

\begin{itemize}
 \item \textbf{Lifted Heston} \citep{abijaber2019affine,el2019perfect}:
  approximates the fractional Riccati--Volterra equation by an
  $n$-dimensional affine jump-diffusion, with $n=5$--$20$ factors
  needed for $H \approx 0.1$. Pricing is via Fourier inversion
  (characteristic function available in closed form), but the
  characteristic function becomes increasingly expensive as $n$ grows.
 \item \textbf{3-factor Bergomi} \citep{bayer2020markovian}:
  parameterises a low-dimensional Markovian system whose joint
  (stock, variance) dynamics approximate the rough Bergomi
  smile. Calibration requires non-convex optimisation over
  $\sim 10$ parameters and the approximation degrades for
  $H < 0.07$, where 3 factors cannot capture the roughness.
\end{itemize}

These Markovian methods share a fundamental limitation: the number of
factors controls the quality of the fractional-kernel approximation, and
the approximation error propagates to option prices in a way that depends
on both the strike and the maturity. In contrast, the spectral surrogate
\emph{directly targets the PSD} of the fGn increments, bypassing the
kernel-approximation step entirely. The implication is that the spectral
surrogate's accuracy is determined by the spectral loss convergence
(Theorem~\ref{thm:main}) rather than by a factor-count parameter, and it
does not degrade as $H \to 0$, a regime where Markovian approximations
require increasingly many factors or break down altogether (Table~\ref{tab:hurst}
shows $|\Delta H| \leq 0.013$ at $H=0.05$, where $n=3$ Markovian
approximations typically incur bias $>0.03$ in recovered $H$).

A secondary advantage is that the spectral surrogate produces paths in
a single forward pass with no matrix exponentials, no Riccati ODE
integration, and no Fourier inversion, it is a pure feed-forward
generator, making it trivially vectorisable and GPU-friendly. The
trade-off is a one-time training cost and the $W_2$ convergence
guarantee that holds only for the Gaussian projection of the law
(Theorem~\ref{thm:main}, Assumption~\ref{ass:gauss}).

\paragraph{Connection to Paper 1.}
The proof structure is a one-assumption extension of \citet{TODO-paper1}:
Bochner (stationary levels) $\to$ Cramér (stationary increments);
all other lemmas carry over unchanged (see the proof-extension guide,
\texttt{proof\_extension\_guide.tex}). The new content is the Hurst
corollary (\cref{cor:hurst}) and the financial validation.

% ═════════════════════════════════════════════════════════════════════════════
\section{Conclusion}
\label{sec:conclusion}
% ═════════════════════════════════════════════════════════════════════════════

We have introduced the Cramér PINN, a neural-network surrogate for the
rough Bergomi and rough Heston models trained to match the exact power
spectral density of fractional Gaussian noise. The main contributions
are: a $W_2$ convergence theorem extending the Bochner PINN framework to
processes with stationary increments; a Hurst exponent consistency
corollary; an empirical demonstration that exact PSD matching achieves
1.21--1.33 vol-point IV surface RMSE under $\rho = 0$, $1.7\times$
better than the asymptotic power-law approximation and near the Monte
Carlo noise floor; a Greek bias analysis showing that the marginal $V$ surrogate carries
an irreducible ATM delta bias of $-0.27$ ($35\%$) when $\rho = -0.9$
under \emph{rough Heston}; a \emph{CausalSpectralFilter} (causal
convolution, same PSD loss) that resolves the circular-FFT wrap-around
pathology, achieving 3.26~vp vs.\ the BLP rBergomi reference (or
\textbf{2.57~vp} with matched-PSD training, leaving only Monte-Carlo
noise) and reducing the rBergomi delta bias to $\lesssim 0.01$; on rough
Heston the same spectral architecture cuts the cross-model bias from
$+0.49$ to $-0.15$ (41$\times$ IV RMSE improvement), and a
Volterra-aware autoregressive surrogate closes the residual to MC noise
(0.80~vp, bias $\leq 0.005$); a Hurst estimator comparison showing that
structure function achieves the smallest bias ($|\text{bias}|\leq 0.011$)
for $H\in[0.05,0.30]$ while PSD slope and R/S are unreliable below $H=0.2$;
a PSD phase diagram establishing $H\approx 0.15$ as the threshold below
which exact (not asymptotic) PSD training is essential; an observation-horizon
sweep recommending $T\leq 2$~yr for IV RMSE below $1.6$~vp; and a live SPX
Hurst estimate $\hat{H}_\mathrm{SPX}=0.145$ from the ATM IV term structure,
recovered by the surrogate to within $|\Delta H|=0.007$.

The CausalSpectralFilter introduced in Section~\ref{sec:greeks} already
handles $\rho \neq 0$ for rBergomi (3.26~vp vs.\ BLP reference).
Remaining natural extensions are: (i) a joint $(V, S)$ surrogate trained
on paired paths, which would close the residual $-0.07$ delta bias
attributable to the V-distribution gap; (ii) the deep hedging application
\citep{buhler2019deep}, for which the causal surrogate is a drop-in path
generator; and (iii) extension to other rough volatility models (RFSV,
multifactor Bergomi) where the same Cramér PSD loss applies.

% ═════════════════════════════════════════════════════════════════════════════
\appendix
\section{Proofs}
% ═════════════════════════════════════════════════════════════════════════════

\subsection{Proof of Theorem~\ref{thm:main}}

We bound the three terms in \eqref{eq:main-bound} separately and then
combine. Throughout, Assumptions~\ref{ass:si}--\ref{ass:gauss} are in force.

\paragraph{Reduction to the Gelbrich formula.}
Under Assumption~\ref{ass:gauss}, both $u_\theta$ and $u_{\mathrm{ref}}$
are centred Gaussian processes; their laws are completely determined by
their covariance operators $C_\theta, C_{\mathrm{ref}}$ on $L^2([0,T])$.
The Bures--Gelbrich formula \citep{gelbrich1990formula} gives
\begin{equation}\label{eq:gelbrich-main}
 W_2^2\!\bigl(\mathcal{L}(u_\theta),\,\mathcal{L}(u_{\mathrm{ref}})\bigr)
 = \bigl\|C_\theta^{1/2} - C_{\mathrm{ref}}^{1/2}\bigr\|_{\HS}^2.
\end{equation}
We bound the right-hand side in four steps.

\paragraph{Step 1 (Square-root Lipschitz).}
Because both processes have stationary increments with Cram\'er spectral
densities satisfying Assumption~\ref{ass:nd}, the covariance operators are
bounded below: $C_\theta, C_{\mathrm{ref}} \geq \mu_{\min} I$ on the
spectral support $[\omega_0,\Omega]$. The operator square-root Lipschitz
bound \citep[Thm.~X.1.1]{bhatia1997matrix} then gives
\begin{equation}\label{eq:sqrt-lip}
 \bigl\|C_\theta^{1/2} - C_{\mathrm{ref}}^{1/2}\bigr\|_{\HS}
 \;\leq\;
 \frac{1}{2\sqrt{\mu_{\min}}}\,
 \bigl\|C_\theta - C_{\mathrm{ref}}\bigr\|_{\HS},
\end{equation}
so that $W_2^2 \leq (4\mu_{\min})^{-1}\|C_\theta - C_{\mathrm{ref}}\|_{\HS}^2$.

\paragraph{Step 2 (Cram\'er kernel to Hilbert--Schmidt norm).}
By Assumption~\ref{ass:si}, the covariance kernel difference is
\[
 \Delta K(t,s) = \int_{\omega_0}^{\Omega}
  (e^{i\omega t}-1)\,\overline{(e^{i\omega s}-1)}\,
  \bigl(S_\theta(\omega)-S_{\mathrm{ref}}(\omega)\bigr)\,d\omega.
\]
Applying Fubini's theorem, $|e^{i\omega t}-1|^2 \leq 4$, and
Cauchy--Schwarz on $[\omega_0,\Omega]$:
\begin{equation}\label{eq:parseval-main}
 \bigl\|C_\theta - C_{\mathrm{ref}}\bigr\|_{\HS}^2
 = \int_0^T\!\!\int_0^T |\Delta K(t,s)|^2\,dt\,ds
 \;\leq\;
 \frac{T^2}{2\pi}\,
 \bigl\|S_\theta - S_{\mathrm{ref}}\bigr\|_{L^2([\omega_0,\Omega])}^2.
\end{equation}
(The factor $2\pi$ absorbs the exact integral
$\int_0^1|e^{i\omega t}-1|^2\,dt = 2-2\sin\omega/\omega \leq 4$.)

Combining Steps~1--2 with \eqref{eq:gelbrich-main}:
\begin{equation}\label{eq:spectral-bound}
 W_2\!\bigl(\mathcal{L}(u_\theta),\,\mathcal{L}(u_{\mathrm{ref}})\bigr)
 \;\leq\;
 \frac{T}{2\sqrt{2\pi\mu_{\min}}}\,
 \bigl\|S_\theta - S_{\mathrm{ref}}\bigr\|_{L^2([\omega_0,\Omega])}
 \;=:\;
 \frac{C_T}{\mu_{\min}^{1/2}}\,
 \bigl\|S_\theta - S_{\mathrm{ref}}\bigr\|_{L^2([\omega_0,\Omega])}.
\end{equation}
This is the spectral term in \eqref{eq:main-bound} with $C_T = T/(2\sqrt{2\pi})$.

\paragraph{Step 3 (Discretisation error).}
The spectral loss $\mathcal{L}(\theta)$ is evaluated on a discrete grid of
$N_s$ frequency points in $[\omega_0,\Omega]$. The surrogate fGn PSD
$S_\theta$ is a piecewise-smooth function inheriting the Sobolev
regularity $H^{H-\varepsilon}$ of fBm increments ($\beta_H =
\tfrac{1}{2}\min(H-\tfrac{1}{2}+\tfrac{d}{2},1)$ for dimension $d$).
Standard FFT quadrature theory gives
\[
 \bigl\|S_\theta^{\mathrm{discrete}} - S_\theta^{\mathrm{continuous}}
 \bigr\|_{L^2([\omega_0,\Omega])}
 \;\leq\; C_{\mathrm{quad}}\,N_s^{-\beta_H},
\]
with $C_{\mathrm{quad}}$ depending on $\|S_\theta\|_{H^{\beta_H}}$
and the geometry of the frequency grid
(cf.\ \citet{TODO-paper1}, Appendix~A, Lemma~A.4).
Inserting into \eqref{eq:spectral-bound} and setting
$C_2 = C_T\,C_{\mathrm{quad}}/\mu_{\min}^{1/2}$ gives the
discretisation term $C_2\,N_s^{-\beta_H}$.

\paragraph{Step 4 (Low-frequency cutoff).}
The spectral window $[\omega_0,\Omega]$ excludes the interval $[0,\omega_0)$
where the fBm PSD $S_H(\omega) = C_H|\omega|^{1-2H}$ diverges for $H < 1/2$.
Restricting to $\omega_0 > 0$ is necessary for Assumption~\ref{ass:nd} to
hold. The penalty for this restriction enters via $\mu_{\min}$: since
\[
 \mu_{\min} \;=\; S_H(\omega_0) \;=\; C_H\,\omega_0^{1-2H},
\]
the constant in \eqref{eq:spectral-bound} satisfies
$\mu_{\min}^{-1/2} = C_H^{-1/2}\,\omega_0^{-(1-2H)/2}$.
For any fixed spectral gap $\|S_\theta-S_{\mathrm{ref}}\|_{L^2}$, the
$W_2$ bound degrades as $\omega_0^{-(1-2H)/2}$ as the window shrinks toward
the origin. Setting
\[
 C_3\,\omega_0^{-(1-2H)/2}
 := \frac{C_T\,\|S_\theta-S_{\mathrm{ref}}\|_{L^2}}{C_H^{1/2}}\,
   \omega_0^{-(1-2H)/2}
\]
identifies the third term; it vanishes as $\omega_0 \to 0$ only when
$H > 1/2$ (the PSD then integrates to zero at the origin).

\paragraph{Combining the three terms.}
The triangle inequality on $W_2$ applied through the spectral estimate at
grid resolution $N_s$ and at frequency cutoff $\omega_0$ gives
\begin{align*}
 W_2\!\bigl(\mathcal{L}(u_\theta),\mathcal{L}(u_{\mathrm{ref}})\bigr)
 &\leq\;
 \underbrace{\frac{C_T}{\mu_{\min}^{1/2}}
 \bigl\|S_\theta-S_{\mathrm{ref}}\bigr\|_{L^2([\omega_0,\Omega])}}_{\text{spectral term}}
 \;+\;
 \underbrace{C_2\,N_s^{-\beta_H}}_{\text{discretisation}}
 \;+\;
 \underbrace{C_3\,\omega_0^{-(1-2H)/2}}_{\text{cutoff}}.
\end{align*}
This is \eqref{eq:main-bound}. The three contributions are independent:
the spectral term vanishes when the trained PSD matches the target;
the discretisation term vanishes as the frequency grid is refined;
the cutoff term reflects the irreducible cost of avoiding the $\omega = 0$
singularity.\hfill$\square$

\subsection{Wood--Chan simulation}
\label{app:wood-chan}

Exact simulation of fGn using circulant embedding
\citep{wood1994simulation, dietrich1997fast}.

The method exploits the fact that the covariance matrix of fGn increments
$(\Delta W^H_1, \ldots, \Delta W^H_N)$ is a Toeplitz matrix
$\Gamma_{jk} = \tfrac{1}{2}(|j-k+1|^{2H} - 2|j-k|^{2H} + |j-k-1|^{2H})$.
Every Toeplitz matrix can be embedded in a circulant $C$ of size $2N$;
if the embedding is positive semidefinite, the square root
$C^{1/2}$ is available via FFT in $O(N\log N)$. A realisation
of fGn is then $\Delta W^H = \mathrm{Re}(\mathcal{F}^{-1}[C^{1/2}z])$
where $z \sim \mathcal{CN}(0, I)$. Davies and Harte (1987) showed the
embedding is always positive semidefinite for $H \in (0,1)$ after
mirroring; \citet{wood1994simulation} and \citet{dietrich1997fast}
provide the practical algorithm. Our implementation is in
\texttt{rbm\_reference.py}, function \texttt{simulate\_fBm\_woodchan}.

% ─────────────────────────────────────────────────────────────────────────────
\bibliographystyle{plainnat}
\bibliography{references}
\end{document}
